
\section{Channels}\fbeg{}\begin{frame}{\ft{The Idea of \q{Channels}}}

\vspace{-.5em}	

{\Large\fontfamily{uhv}\selectfont

\hspace*{20pt}\begin{minipage}{\textwidth}
\vspace{4pt}
		
\begin{lightquadblockc}{0,0.1,0.1}{\parbox{21cm}{\centering \vspace{3pt}}}
\begin{center}\begin{minipage}{1.05\textwidth}
{\Huge \setlength{\leftmargini}{30pt}\begin{enumerate}
\dmitem {\dsp} Group together input and/or output parameters with related (maybe special-purpose) semantic profiles 
\vspace{20pt}
\dmitem {\dsp} Language extensions can be implemented via special types of channels 
(e.g., supporting \textsl{queries} inside \textsl{scripting} formats by 
building channels that model query parameters) 
\vspace{20pt}
\dmitem {\dsp} Back-door channels can implement proof-checks and other DbC features
\vspace{20pt}
\dmitem {\dsp} Channels provide convenient search-points for static analysis 
or runtime monitors
\vspace{20pt}
%\footnotemark[2]\footnotetext{\Large \raisebox{7pt}{\em{2}} 
%	\hspace{7pt}See Slides 11 and 12 for details.}
\end{enumerate}
}\end{minipage}
\end{center}
\end{lightquadblockc}
\end{minipage}

}

\end{frame}

